% Part: normal-modal-logic
% Chapter: syntax-and-semantics
% Section: introduction

\documentclass[../../../include/open-logic-section]{subfiles}

\begin{document}

\olfileid{nml}{syn}{int}

\olsection{جان پہچان}

امکانی منطق \emph{امکان تے ضرورت ورگی جہت والے قضیاں} تے
اوہناں دے آپس دے منطقی نتیجیاں بارے اے۔ ایہو جیہے قضیاں دیاں
کجھ مثالاں ہِٹھ نیں۔
\begin{enumerate}
\item ضروری اے کہ $2+2=4$۔
\item ایہ لازماً ممکن اے کہ بھلکے مینہ پوے۔
\item جے~$!A$ دا ہونا لازماً ممکن اے، تاں~$!A$ دا ہونا ممکن اے۔
\end{enumerate}
امکان تے ضرورت توں علاوہ ہور جہتاں وی نیں۔ ہور اک-دلیلی
رابطے وی جہتی رابطے منّے جاندے نیں، جیویں «~$!A$ ہونا چاہیدا
اے»، «~$!A$ ہووے گا»، «ڈینا جاندا اے کہ~$!A$»، یا
«ڈینا منّدا اے کہ~$!A$»۔

امکانی منطق پہلاں ارسطو دی \emph{De
  Interpretatione} وچ سامنے آؤندی اے۔ اوہنے سبھ توں پہلاں
دھیان دِتّا کہ ضرورت توں امکان نکلدا اے، پر اُلٹی گل نہیں۔
امکان تے ضرورت نوں اک دوجے دے ذریعے متعین کیتا جا سکدا اے۔
جے $!A \land !B$ دا سچ ہونا ممکن اے، تاں $!A$ دا سچ ہونا
ممکن اے تے $!B$ دا سچ ہونا وی ممکن اے، پر اُلٹی گل نہیں۔
تے جے $!A \to !B$ ضروری اے، تاں جے $!A$ ضروری اے،
~$!B$ وی ضروری اے۔

امکانی منطق دا پہلا جدید طریقہ سی۔ آئی۔ لوئس دے کم توں
آیا، جیہڑا لوئس تے لینگفورڈ دی کتاب \emph{Symbolic Logic}
(1932) تک پُہنچیا۔ لوئس تے لینگفورڈ، لزوم نوں مادی شرطیے
نال ظاہر کرن توں مطمئن نہیں سن۔ $!A \lif !B$، «$!A$ توں
$!B$ لازم اے» دی تھاں چنگی طرح کم نہیں کردا۔ ایس لئی اوہناں
لزوم نوں «ضروری اے کہ جے $!A$، تاں $!B$» وجوں متعین
کرن دی تجویز دِتّی، تے ایہنوں $!A \strictif !B$ لکھیا۔
وکھ وکھ خاصیتاں سمجھدیاں لوئس نے پنج امکانی نظام پچھانے،
\Log{S1}، \ldots، \Log{S4}، \Log{S5}۔ آخری دو اجے وی
ورتے جاندے نیں۔

لوئس تے لینگفورڈ دا طریقہ سراسر \emph{نحوی} سی۔ اوہناں
موزوں مسلّمے تے قاعدے چُنے تے پرکھیا کہ اوہناں نال کیہ
ثابت ہو سکدا اے۔ معنیات دا طریقہ لمّے چِر تک ہتھ نہ آیا۔ پہلی
کوشش روڈولف کارناپ نے \emph{Meaning and Necessity} (1947)
وچ \emph{حالت دے بیان} دا خیال ورت کے کیتی۔ ایہ ایٹمی
جملیاں دا اکٹھ اے، یعنی اوہ جملے جیہڑے اوس حالت دے بیان
وچ «سچ» نیں۔ سچ ہون دی تعریف ہر طرح دے جملے $!A$ تک
ودھا کے، کارناپ $!A$ نوں \emph{لازماً سچ} منّدا اے جے اوہ
حالتاں دے سارے بیاناں وچ سچ ہووے۔ کارناپ دا طریقہ \emph{وار
وار لگیاں} جہتاں نوں نہیں سنبھال سکدا سی۔ «ممکن اے کہ
ضروری اے \ldots ممکن اے کہ $!A$» ورگے جملے ہمیشہ
سبھ توں اندرلی جہت تک گھٹ جاندے نیں۔

امکانی معنیات وچ وڈی پیش رفت ساؤل کرپکے دے مضمون
«A Completeness Theorem in Modal Logic» (JSL 1959) نال
ہوئی۔ کرپکے دا کم لائبنز دے ایس خیال اُتے ٹکیا سی کہ جملہ
لازماً سچ اے جے اوہ «سارے ممکن جہاناں وچ» سچ اے۔ پر ایس
خیال وچ وی کارناپ والی دِقّت اے۔ کسے جہان $w$ یا حالت دے
بیان $s$ وچ جملے دا سچ ہونا $w$ اُتے بالکل منحصر نہیں۔
سو کرپکے نے منّیا کہ جہان \emph{رسائی دے تعلق} $R$ نال
جُڑے نیں۔ «ضروری اے کہ $!A$» دی شکل والا جملہ جہان $w$
وچ اُدوں تے صرف اُدوں سچ اے جدوں $!A$ سارے اوہناں جہاناں
$w'$ وچ سچ اے جنہاں تک \emph{رسائی} $w$ توں اے۔
ایس طریقے دا کوئی روپ دین والی معنیات نوں کرپکے معنیات
آکھدے آں۔ ایہناں نال کئی امکانی منطقاں تیزی نال اگّے ودھیاں۔

کرپکے معنیات مطابق سمجھیے، تاں امکانی منطق سانوں دسدِی
اے کہ \emph{تعلقی ساختاں} «اندروں» کیویں دِسدیاں نیں۔
تعلقی ساخت بس اوہ سیٹ اے جیہدے اُتے دو رکنی تعلق اے۔
مثلاً جماعت دے طالب علماں دا سیٹ، جیہڑا اوہناں دے سماجی
تحفظ دے نمبراں مطابق ترتیب دِتّا ہووے، اک تعلقی ساخت اے۔
تعلقی ساختاں کئی طرح دے میدان وچ ہوندیاں نیں۔ دنیا دیاں
حالتاں دے اک دوجی دے لحاظ نال ممکن ہون توں علاوہ، کسے
کارندے دیاں علمی حالتاں علمی امکان نال جُڑ سکدیاں نیں۔
ایسے طرح بدل دے نظام دیاں حالتاں اوہناں دیاں حالتی
تبدیلیاں نال جُڑ سکدیاں نیں۔ امکانی منطق نال ایہناں
ساریاں دا نمونہ بنا سکدے آں۔ پہلی صورت توں صدق نال جُڑی
عام امکانی منطق ملدی اے۔ ہور صورتاں توں علمی منطق، عمل تے
تبدیلی دی منطق وغیرہ ملدیاں نیں۔

ساڈی توجہ اک خاص پاسے اُتے اے جیہنوں امکانی منطق دے
ماہر «مطابقت دا نظریہ» آکھدے نیں۔ کرپکے دیاں اہم ابتدائی
لبھتاں وچوں اک ایہ اے کہ رسائی دے تعلق~$R$ دیاں کئی
خاصیتاں، جیویں ایہدا تعدّی والا یا متقابل ہونا، مناسب
«امکانی سانچیاں» نال \emph{امکانی زبان دے اندر ای} پچھانیاں
جا سکدیاں نیں۔ مثلاً امکانی منطق دے ماہر آکھدے نیں کہ
$R$ دا انعکاسی ہونا سانچے «جے لازماً~$!A$، تاں~$!A$»
نال «مطابق» اے۔ اسیں زیادہ تر امکانی منطق دے کجھ کلاسیکی
نظاماں، مثلاً \Log{S4} تے \Log{S5}، دا مطابقتی نظریہ
ویکھاں گے۔ ایہ نظام سانچیاں \Ax{D}، \Ax{T}، \Ax{B}،
\Ax{4} تے~\Ax{5} دیاں ترکیباں توں ملدے نیں۔

\end{document}
